\section{Eenvoudige thermodynamische systemen}
\subsection{PV-diagrammen (zuivere stof)}

\begin{definition}
    \textbf{Vloeistofverzadigingslijn}: De lijn in een PV-diagram waarbij vloeistof en damp in evenwicht zijn. Deze lijn scheidt het gebied van de vloeistoftoestanden van het gebied van de damp-toestanden.
\end{definition}
\medskip
\begin{definition}
    \textbf{Kritische temperatuur}: De temperatuur waarbij de vloeistofverzadigingslijn eindigt. Boven deze temperatuur kunnen vloeistof en damp niet meer in evenwicht zijn, en is er geen onderscheid meer tussen de twee fasen.
\end{definition}
\medskip
\begin{definition}
    \textbf{Kritische isotherm}: De isotherm die door het kritische punt gaat. Deze isotherm heeft een horizontaal inflectiepunt bij het kritische punt, wat betekent dat de eerste en tweede afgeleide van de druk naar het volume gelijk zijn aan nul op dat punt.
\end{definition}
\medskip
\begin{definition}
    \textbf{Kritisch punt K}: Het punt in het PV-diagram waar de vloeistofverzadigingslijn eindigt. Op dit punt zijn de eigenschappen van de vloeistof en de damp identiek, en er is geen onderscheid meer tussen de twee fasen.
\end{definition}
\medskip
\begin{definition}
    \textbf{Kritische druk}: De druk die overeenkomt met het kritische punt in het PV-diagram. Dit is de druk waarbij de vloeistofverzadigingslijn eindigt en de kritische isotherm een horizontaal inflectiepunt heeft.
\end{definition}
\medskip
\begin{definition}
    \textbf{Verzadigde vloeistof}: Een vloeistof die in evenwicht is met haar eigen damp. Deze vloeistof heeft een constante druk en temperatuur.
\end{definition}
\medskip
\begin{definition}
    \textbf{Verzadigde damp}: Een damp die in evenwicht is met zijn eigen vloeistof. Deze damp heeft een constante druk en temperatuur.
\end{definition}
\medskip
\begin{definition}
    \textbf{dampverzadigingslijn}: De lijn in een PV-diagram waarbij damp en vloeistof in evenwicht zijn. Deze lijn scheidt het gebied van de damp-toestanden van het gebied van de vloeistoftoestanden.
\end{definition}
\medskip
De dampverzadigingslijn en de vloeistofverzadigingslijn ontmoeten elkaar bij het kritische punt K, waar de eigenschappen van de vloeistof en de damp identiek zijn.

\begin{definition}
    \textbf{Tripellijn}: De lijn in een PV-diagram waarbij de vloeistof-damp en vaste-stof-damp in evenwicht zijn.
\end{definition}
\medskip
\begin{definition}
    \textbf{Tripelpunt}: Het punt in een PV-diagram waar de drie fasen van een stof (vloeistof, damp, en vaste stof) in evenwicht zijn. Op dit punt kunnen alle drie de fasen tegelijkertijd bestaan.
\end{definition}
\medskip
\begin{definition}
    \textbf{Sublimatielijn}: De lijn in een PV-diagram waarbij de vaste stof en de damp in evenwicht zijn. Deze lijn scheidt het gebied van de vaste-stof-toestanden van het gebied van de damp-toestanden.
\end{definition}
\medskip
